\begin{figure}[H]
\centering
\begin{tikzpicture}[font=\small,>=Latex,
  caixa/.style={draw=blue!65!black,fill=blue!5,text width=3.5cm,minimum height=2.35cm,align=center}]
\node[caixa] (h) at (0,0) {\textbf{Host / CPU}\\[4pt]Vetores em \texttt{malloc}\\$a$, $b$, $c$ e referência};
\node[caixa,fill=orange!12,draw=orange!70!black] (d) at (8,0)
{\textbf{Dispositivo / GPU}\\[4pt]Dados acessíveis no dispositivo\\$c[i]=a[i]+b[i]$};
\draw[->,thick] (h.30)--node[above,align=center,font=\footnotesize]
{Entradas: $a$ e $b$\\80 MB de dados} (d.150);
\draw[->,thick] (d.210)--node[below,align=center,font=\footnotesize]
{Saída: $c$\\40 MB de dados} (h.330);
\node[align=center,text width=12.3cm] at (4,-2.65)
{Cada vetor possui dez milhões de elementos.\\A referência fica no host para a validação.};
\end{tikzpicture}
\caption{Modelo de movimentação dos dados úteis entre host e dispositivo. O desenho representa as entradas e a saída necessárias à operação, sem medir o tráfego efetivo do runtime.}
\label{fig:offload}
\end{figure}
